% RESEARCH_DRAFT_SHA256: 7daaa10a8ef705a5665e92eb4a47fb68b9fa3c248567ba52fb6ca0dab6e81ea8
% Layout and source-reference wrappers expanded for the web renderer.
\documentclass[11pt]{article}
\usepackage[margin=25mm]{geometry}
\usepackage{amsthm}
\usepackage{fontspec,unicode-math}
\setmainfont{Latin Modern Roman}
\setmathfont{Latin Modern Math}
\usepackage{amsmath,amssymb,mathtools,xurl,hyperref}
\hypersetup{colorlinks=true,urlcolor=blue,linkcolor=blue}
\setcounter{secnumdepth}{0}
\setlength{\parindent}{0pt}
\setlength{\parskip}{5pt}
\setlength{\emergencystretch}{3em}
\newcommand{\sref}[2]{\hyperlink{source:#1:#2}{[S#2]}}
\newcommand{\cataloguescope}[1]{\paragraph{Source record.}#1}
\newcommand{\floor}[1]{\left\lfloor#1\right\rfloor}
\newtheorem{proposition}{Proposition}
\newtheorem{lemma}[proposition]{Lemma}
\begin{document}
\paragraph{Research manuscript.} OpenAI \texttt{gpt-6-astra}, reasoning \texttt{ultra}; 11 October 2026.
\paragraph{Preserved source body.} The complete catalogue mathematical text and sources follow with layout and source links expanded, including source qualifications and its historical preparation notes. New research begins at the section headed Research attempt.

% UnsolvedMath ID 20002130; source catalogue code AIM-INFRASTRUCTURE-0030
\section{A versioned localization and transport certificate}\label{20002130}
{\small UnsolvedMath ID 20002130 · Computer Science}
\subsection{Definitions and mathematical statement}

\paragraph{Shared notation.}Standard number systems are $\mathbb N,\mathbb Z,\mathbb Q,\mathbb R,\mathbb C$. Specialized notation and index ranges follow the cited source and the definitions below.


An international educational deployment introduces a resource or teaching approach in another national setting. Localization concerns language and context; institutional adaptation concerns curricula, infrastructure, assessment, and available support. A cross-context comparison needs comparable interpretations of the measured outcomes.
\paragraph{Record scope and qualifications.}
The catalogue title is retained as metadata. Named certificates, contracts, bounds or frameworks in that title are not asserted as results of the original source extract.
Neither the resource nor the destination contexts, constraints, or effectiveness measures are stated.

Catalogue screening limitation: Educational assessment, adoption, or study-design discussion; the supplied question does not state a mathematical target with defined assumptions and outcome.

The source snapshot labels this record partially solved; this is catalogue metadata rather than independent verification of a mathematical result or the current completion of a project.
\paragraph{Statement recorded in the supplied source.}
32. Expanding to other countries \href{http://aimath.org/pastworkshops/oerassessmentproblems.pdf}{[S2]}
\paragraph{Source status.}
The supplied catalogue labels this source record partially solved; that metadata was checked against the supplied snapshot on 2026-10-11. The label is preserved for provenance and is not an independent verification of present open status or mathematical correctness.
\subsection{Short English statement}
Consider extending the educational tools or practices to additional countries.
\subsection{Sources}
\begin{itemize}
\item[S1] Ulam AI and the UnsolvedMath Contributors, UnsolvedMath: A Curated Collection of Open Mathematics Problems, record 20002130 (AIM-INFRASTRUCTURE-0030). Catalogue curation and metadata: CC BY 4.0. The source extract is retained; definitions, English summaries and qualifications are editorial additions.
\url{https://huggingface.co/datasets/ulamai/UnsolvedMath}.
\item[S2] AIM Workshop Problem Lists, source linked by catalogue record AIM-INFRASTRUCTURE-0030.
\url{http://aimath.org/pastworkshops/oerassessmentproblems.pdf}.
\end{itemize}
\paragraph{Preparation and review.}
The original source and UnsolvedMath contributors retain attribution for the source extract. Codex (GPT-6; exact serving revision not exposed) assisted with definitions, English summaries, scope notes and formatting. No new human mathematical authorship or independently verified solution is asserted. Permanent source ID 20002130 and code AIM-INFRASTRUCTURE-0030 checked against the complete supplied registry and current repository. Distinct upstream records, including repeated questions and broad or damaged extracts, are retained for complete coverage of the source CS category. A repeated record is not a new independent conjecture. Complete mathematical-equivalence and current-literature review has not been performed. Preparation and final review follow contribution prompts C and D, with source gaps disclosed. Structural and rendering checks do not establish a new theorem or open status.
\label{end:20002130}

\section{Research attempt}
\subsection{Literature, source context and scope}
The primary AIM notes describe a brainstorming session. The original excerpt and all catalogue qualifications remain above. They do not specify a mathematical theorem or a complete empirical study. The formal model below is an explicitly limited research route introduced for this attempt, not a reconstruction of a missing original conjecture. The supplied catalogue's earlier research summary was inspected as context; its named framework is not adopted as an established solution. The elementary result below is not claimed novel, and no educational effectiveness data were collected.

\subsection{Attempt: distinguishing localization effects from learner differences}
For the source's international-expansion agenda, I tested a specific measurement ambiguity. Declare the population response model
\[
\operatorname{logit}P(Y_{ij\ell}=1)=\theta_i-b_j-\delta_\ell,
\qquad \operatorname{logit}(p)=\log\frac p{1-p},
\]
where $\theta_i$ is learner ability, $b_j$ item difficulty and $\delta_\ell$ a language-version shift. This Rasch-style additive model is an assumption, not an established cross-country measurement law.
\begin{proposition}
If each learner is observed in only one language, a language group's ability shift is indistinguishable from an equal shift in that language's item-version parameter. If the same learner answers the same item in two languages under this model, the difference of the two population logits identifies the difference of language shifts.
\end{proposition}
\begin{proof}
For all learners observed only in language $\ell$, replace $\theta_i$ by $\theta_i+c$ and $\delta_\ell$ by $\delta_\ell+c$. Every observed linear predictor in that group is unchanged; all other groups are untouched. Hence every observed response probability is identical while the group ability levels and language parameter differ. No likelihood based on those probabilities distinguishes the two descriptions.

For a fixed learner $i$ and item $j$ appearing in languages $\ell$ and $m$, subtract the model equations. The common $\theta_i-b_j$ cancels, leaving
\[
\operatorname{logit}P(Y_{ij\ell}=1)-\operatorname{logit}P(Y_{ijm}=1)
=-(\delta_\ell-\delta_m).
\]
Thus the population logit difference identifies that contrast under the stated common-ability and common-item model.
\end{proof}
\subsection{Verification and remaining gaps}
The predictor $0.7-0.2-0.4=0.1$ is unchanged after replacing ability 0.7 by 1.2 and language shift 0.4 by 0.9. This exact identity checks the confounding construction. The bilingual formula uses population probabilities, not the logit of a single binary observation; finite estimation, carryover and translation-specific skill demands require separate work. No national-group ability comparison or transfer of educational effectiveness is inferred. The failed route is interpreting cross-language score differences without a valid measurement anchor.

\section{References and source attribution}
American Institute of Mathematics, workshop discussion notes on open educational resources and assessment, \url{https://aimath.org/pastworkshops/oerassessmentproblems.pdf}. The 16-page primary PDF was opened on 11 October 2026 and its numbered brainstorming item was checked against the supplied record. The source supports the provenance and agenda, not the additional modeling assumptions or an empirical educational conclusion. Catalogue and contributor attribution are retained in the reproduced Sources section. The worked mathematics is proved here without reliance on an unverified external theorem.


\paragraph{Integration review.} Codex (GPT-6; exact serving revision not exposed) reviewed the written arguments and expanded layout and source-reference wrappers. Repository metadata and the production web renderer were checked using Python 3.13 and JavaScript with MathJax 3.2.2. This model-assisted review is not independent human proof verification; no PDF compilation or formal proof checking is asserted.

\section{Completion Estimate}
\noindent\textbf{COMPLETION: 0\%}
No substantiated new progress toward the entire original target is claimed. The worked results and failed route are delimited above; a known result, illustrative model, or auxiliary obstruction is not counted as original-target completion.

\section{Final status and attribution}
\textbf{Attempt status: unresolved.} This describes this research attempt, not a new ruling on the source catalogue's status or on external literature claims.
The mathematics and this self-audit were generated by OpenAI \texttt{gpt-6-astra}, reasoning \texttt{ultra}, on 11 October 2026. Codex (GPT-6; exact serving revision not exposed) orchestrated the batch and contributed formatting/repository review. No human mathematical authorship, independent human proof checking, or formal verification is claimed. The historical preparation paragraphs in the reproduced catalogue body concern that body only.

Contribution prompt A was applied to the complete definition and supplied source record, followed by the adversarial checks of prompt B. Local reproduction files are \texttt{batch008/checks.py}, \texttt{computations.json}, and the per-ID audit JSON. They were not public links at preparation time. Numerical checks are finite and do not prove unrestricted results. No LaTeX engine was found locally; compilation was not performed by this research agent.

\end{document}
